% !TEX root = ../thesis_book_0421052099.tex
\chapter{Background and Related Work}\label{ch:background}

This chapter lays out the concepts and notation that the technical chapters rely
on. It also surveys the work that AGR builds on and competes with. The chapter
is kept short on purpose. It covers only material that is actually used later.
Standard results are stated, not derived. The background comes first: knowledge
graphs, question answering over them, language models as reasoning engines,
retrieval augmentation, and search. The survey follows. It moves from static
graph-augmented generation, through path retrieval and agentic exploration, to
verification and factuality measurement. The chapter closes with a structured
comparison of the seven closest systems. That comparison does two things. It
points to the gap this thesis addresses. And it justifies the choice of
baselines in \Cref{sec:baselines}.

\section{Knowledge Graphs}\label{sec:kgs}

\subsection{Triples, Entities, and Relations}

A \emph{knowledge graph} $G = (E, R, T)$ consists of a set of entities $E$, a
set\index{triple}\index{knowledge graph} of relation types $R$, and a set of
triples $T \subseteq E \times R \times E$. A triple\index{triple} $(h, r, t)$
asserts that relation $r$ holds from head entity $h$ to tail entity $t$ --- for
example, \textsf{(Ada Lovelace, place\_of\_birth, London)}. Relations are
directed and typed. Both directions carry meaning: the inverse of
\textsf{place\_of\_birth} is a legitimate but distinct relation.

Two properties matter for what follows. Triples are \emph{atomic}. Each asserts
exactly one fact. So a claim can be checked against the graph by testing for the
existence of specific edges. They are also \emph{composable}: a chain $(e_0,
r_1, e_1), (e_1, r_2, e_2), \dots$ expresses a multi-step relationship that no
single triple states. That composability is what makes a knowledge graph the
natural substrate for multi-hop questions.

\subsection{Freebase: MIDs, Schema, and Mediator Nodes}

Freebase~\cite{bollacker2008freebase} is a large, collaboratively-built
knowledge\index{Freebase} base, frozen since 2016. It is also the substrate for
the standard knowledge graph question answering benchmarks. Three of its
characteristics affect any system built on it. Entities are identified by opaque
\emph{machine identifiers}\index{MID} (MIDs) such as \textsf{m.02mjmr}, with
human-readable names attached as a separate property. Names are neither unique
nor stable. So entity identity and entity surface form must be handled
separately. Relations are namespaced paths such as
\textsf{people.person.place\_of\_birth}, encoding domain, type, and property.
They are drawn from a vocabulary of tens of thousands that includes many
carrying no factual content, such as internal type and key predicates.

Most consequentially, Freebase represents $n$-ary facts through
\emph{mediator}\index{mediator node} or \emph{compound value type} (CVT) nodes.
A fact with more than two participants --- an actor playing a particular
character in a particular film --- is not a single edge but an anonymous
intermediate node linking the participants. So a relationship a user would
describe as one hop is two edges in the graph. Any system reporting hop counts
over Freebase must state how it treats these nodes. The choice changes every
depth measurement. Systems that ignore them lose the facts they
encode~\cite{xu2024generate}.

\subsection{Property Graphs, Neo4j, and Cypher}

The triple model above is the RDF view. The alternative is the \emph{property
graph}\index{property graph} model, in which nodes and edges are both typed and
may carry arbitrary key--value properties. Neo4j~\cite{robinson2015graph}
implements it. Cypher~\cite{francis2018cypher} queries it, with a \textsf{MATCH}
clause that expresses graph patterns directly. A triple store maps onto it
naturally --- entities become nodes carrying their MID and name, relations
become typed edges. Traversal patterns that would require recursive joins in a
relational schema are expressed as literal path patterns instead. That is the
practical reason for choosing a graph store here.

\subsection{Knowledge Graph Question Answering}\label{sec:kgqa}

\paragraph{Multi-Hop Questions and Constraint Satisfaction}

Knowledge graph question answering (KGQA)\index{KGQA} takes a
natural-language\index{multi-hop question} question and a graph and returns the
entity or entities that answer it. A question is \emph{$k$-hop} if answering it
requires traversing $k$ edges from the entities named in the question --- its
\emph{topic entities} --- to the answer. Beyond hop count, questions vary in
\emph{constraints}: conditions that filter candidates rather than extending the
chain. ``Which of Tolstoy's novels was published after 1875?'' is one hop plus a
temporal constraint. A correct answer has to satisfy both at once. A system that
walks the hop correctly but quietly drops the constraint returns every novel,
not just the ones published after 1875. The result is a superset of the answer.
Hop-count analysis alone never exposes that failure.

Two families of approach have historically divided this field. \emph{Semantic
parsing} methods translate the question into a formal query and execute it. That
approach is exact and interpretable, but it requires the parse to be exactly
right~\cite{berant2013semantic,yih2016value}. \emph{Information retrieval}
methods instead pull out a relevant subgraph and rank candidate answers inside
it. They fail more gently, but they are harder to see into. The agentic systems
considered here sit closer to the second family. But they get much of the first
family's transparency back by logging the traversal itself.

\paragraph{Evaluation Conventions: Hits@1 and F1}

% makeindex reads @ as the sort/actual separator, so an unquoted Hits@1 prints
% as a bare "1" filed under H. The " quotes the next character.
Two metrics are standard. \emph{Hits@1}\index{Hits"@1} scores a question
$1$\index{F1} if any entity the system asserts is in the gold answer set, and
$0$ otherwise. It is a simple right-or-wrong score per question. It is also the
headline number in most of this literature. The name is historical and slightly
misleading. It was coined for systems that returned a \emph{ranked} list, where
only the top entry counted. It has since been carried over to systems that
return an unordered set. There it becomes the set-intersection test just stated.
This thesis uses the set-intersection form throughout, because that is what the
systems it compares against report. \Cref{sec:metrics} pins down the
string-matching rule exactly.

Many questions have more than one correct answer. Hits@1 alone rewards a system
that names one correct entity and ignores the rest. \emph{F1} fills that gap. It
is computed per question between the predicted answer set $A$ and the gold set
$A^{*}$: \[ P = \frac{|A \cap A^{*}|}{|A|}, \qquad R = \frac{|A \cap
A^{*}|}{|A^{*}|}, \qquad F_1 = \frac{2PR}{P + R}, \] and then averaged over
questions. The difference between $P$ and $R$ matters a great deal in this
thesis. A system that asserts many entities can score high recall while its
precision falls. \Cref{sec:limitations} argues that exactly this failure goes
unmeasured when only Hits@1 is reported.

\section{Large Language Models as Reasoning Engines}\label{sec:llms}

\subsection{In-Context Learning, Chain-of-Thought, and the ReAct Loop}

Sufficiently large language models perform tasks specified entirely in their
prompt, without gradient updates --- \emph{in-context learning}\index{in-context
learning}~\cite{brown2020language}. Prompting the model to produce intermediate
reasoning steps before its final answer, \emph{chain-of-thought}
prompting~\cite{wei2022chain}, substantially improves performance on multi-step
problems. Both are used here: the agent's planner, scorer, evaluator, and
verifier are prompted components of a single frozen model. No component is
fine-tuned. That matters for comparability, since several systems surveyed in
\Cref{sec:comparison-table} are fine-tuned and therefore not directly comparable
to prompting methods.

A language model can also be given \emph{tools}: named operations it may invoke,
whose results re-enter its context~\cite{schick2023toolformer}. The
\emph{ReAct}\index{ReAct} pattern~\cite{yao2023react} interleaves reasoning and
action in a loop --- the model reasons about what to do, acts, observes the
result, and reasons again. That loop is the control structure underlying agentic
graph navigation. Where the tools are graph operations, each iteration extends a
traversal. The sequence of tool calls is a complete record of how the answer was
reached.

\subsection{A Working Definition of Hallucination for This
Thesis}\label{sec:hallucination-def}

``Hallucination'' means different things in different papers. This thesis adopts
a narrow,\index{hallucination} operational definition tied to the knowledge
environment: an asserted entity is \emph{ungrounded} if it does not exist in the
graph, or exists but is not reachable from the question's topic entities within
the traversal budget. A claim is \emph{unsupported} if the traversed triples do
not entail it.

Two consequences of this definition should be stated plainly. The definition is
relative to the environment: an entity absent from the graph is ungrounded even
if true in the world. And groundedness is not correctness --- an answer can be
perfectly grounded, every entity real and reachable, and still be the wrong
answer to the question. \Cref{sec:results} shows that this gap is not
hypothetical. \Cref{sec:limitations} is built on it.

\subsection{Retrieval-Augmented Generation}\label{sec:rag}

Dense retrieval encodes queries and documents into a shared vector space and
retrieves by nearest neighbour, capturing semantic similarity where lexical
matching fails~\cite{karpukhin2020dense}. Sentence-level encoders such as
Sentence-BERT~\cite{reimers2019sentence} produce embeddings suitable for direct
cosine comparison. In this thesis embeddings serve two bounded purposes ---
linking question phrases to graph entities, and pre-ranking relation candidates
during exploration. They never establish that a fact is true. That job belongs
to the symbolic triples alone. Graph-based retrieval swaps flat passages for
structured context~\cite{edge2024local,gutierrez2024hipporag}. What gets
retrieved is a subgraph rather than a document. So the relationships between
facts come along explicitly. The distinction that matters here is not graph
versus vector. It is \emph{static versus interleaved}: does retrieval happen
once, before reasoning begins, or again and again as reasoning proceeds?

\subsection{Search over Graphs}\label{sec:search}

\emph{Best-first search}\index{best-first search} keeps a frontier of candidate
nodes. A heuristic scores each one by how promising it looks, and the search
expands the most promising one first~\cite{hart1968formal,pearl1984heuristics}.
\emph{Beam search} trims the frontier to the $b$ highest-scoring candidates at
each depth. That gives up completeness in exchange for a fixed memory and
computation budget. With an expensive scoring function --- here, a language
model call --- the beam width is also the principal cost control. Bounded search
commits to choices that may prove wrong. \emph{Backtracking}\index{backtracking}
recovers by abandoning the current branch and resuming from a previously
unexplored alternative. That recovery requires maintaining a stack of such
alternatives and a trigger condition for when to pop it. Both are specified in
\Cref{sec:evaluator}.

\paragraph{Relation to Monte Carlo Tree Search: A Terminological
Clarification}\label{sec:mcts-clarification}

The original proposal for this work described its navigation strategy
as\index{best-first search} ``inspired by Monte Carlo Tree Search''. That
description is withdrawn here. The reason is worth stating explicitly. Monte
Carlo Tree Search~\cite{kocsis2006bandit,browne2012survey} is defined by four
phases: selection, expansion, \emph{simulation} (a rollout to a terminal state),
and \emph{backpropagation} of the rollout's value to the ancestors of the
expanded node. The two phases that give MCTS its statistical character are
simulation and backpropagation. Together, they let node values be estimated from
sampled outcomes rather than from a fixed heuristic.

The method used in this thesis has neither. There are no rollouts. No value is
propagated back up the search tree: each candidate is scored once, by a fixed
function combining embedding similarity with a language-model judgement. That
score is never revised in light of what is found later. What the method is, is
\emph{guided best-first search with a bounded beam and backtracking}. It is
described that way throughout. The contrast is not hypothetical. Genuine tree
search has been applied to this exact task: KBQA-o1~\cite{luo2025kbqao1} drives
an agentic knowledge base question answering process with full MCTS --- policy
and reward models, rollouts, and value backpropagation --- over Freebase.
Reasoning with Trees~\cite{shen2025reasoning} likewise builds an explicit search
tree. That those systems exist is precisely why the terminology matters here:
describing a single-pass scoring heuristic as MCTS-inspired would invite
comparison with machinery this framework does not implement.

\subsection{Statistical Tools Used in This Thesis}\label{sec:stats}

Three tools recur, each chosen for a specific property of the
data.\index{McNemar}\index{bootstrap}\index{Cohen's kappa}

\emph{Bootstrap confidence intervals}~\cite{efron1979bootstrap} put a number on
sampling uncertainty. The evaluated questions are resampled with replacement,
and the metric is recomputed on each resample. The $2.5$th and $97.5$th
percentiles of the resulting distribution give a 95\% interval. That is the
\emph{percentile} method. It is used here rather than the bias-corrected
variants. The intervals are read as a sense of scale for a reported difference,
not as a hypothesis test. They also make no distributional assumption about
metrics such as per-question F1.

\emph{McNemar's test}~\cite{mcnemar1947note} asks whether two systems differ in
accuracy when both are run on the \emph{same} questions. It looks only at the
discordant pairs --- questions that one system gets right and the other gets
wrong --- and tests whether the split between the two discordant counts is more
lopsided than chance would produce. Paired testing is essential here, because
every system is run on the same question samples. An unpaired test would throw
the pairing away and lose power.

\emph{Cohen's $\kappa$}~\cite{cohen1960coefficient} measures agreement between
two annotators on a categorical judgement, correcting for agreement expected by
chance: \[ \kappa = \frac{p_o - p_e}{1 - p_e}, \] where $p_o$ is observed
agreement and $p_e$ is the agreement expected if both annotators labelled
independently at their observed marginal rates. It is used to validate the
automatic groundedness judge against human annotation. Conventional
interpretation follows Landis and Koch~\cite{landis1977measurement}, on whose
scale values between $0.61$ and $0.80$ denote substantial agreement.
\section{Static Graph-Augmented Generation}\label{sec:static-graphrag}

\subsection{GraphRAG and Query-Focused Summarization}

GraphRAG~\cite{edge2024local} constructs a knowledge graph from a text
corpus\index{GraphRAG} using a language model. It partitions the graph into
communities and pre-computes summaries at each level. So a query can be answered
from structured context rather than retrieved passages. Its strength is global
questions --- those whose answer depends on the corpus as a whole rather than
any single passage. Flat retrieval handles these poorly.

The relevant limitation is that retrieval remains a single step preceding
generation. The context is assembled before the model has begun reasoning. So it
cannot be conditioned on intermediate findings. GraphRAG also builds its graph
by language-model extraction. This is appropriate for unstructured corpora, but
it introduces extraction error into the reference itself --- a consideration
that shapes the environment design in \Cref{sec:env-goals}.

\subsection{HippoRAG and Memory-Structured Retrieval}

HippoRAG~\cite{gutierrez2024hipporag} draws on the hippocampal indexing theory
of human memory. It combines a knowledge graph with a personalised PageRank
retrieval step, so that a single query can integrate evidence distributed across
many passages. It substantially improves multi-hop retrieval over flat baselines
at low cost.

Its retrieval is nonetheless a single graph-propagation step, not an
agent-directed traversal: the propagation is fixed. The model does not choose
where to look next. Temporal extensions such as T-GRAG~\cite{li2025tgrag}
address a different axis --- conflicting and redundant facts across time. This
is orthogonal to the navigation question. So it is noted here for completeness
rather than compared against.

\subsection{Path-Retrieval and Semantic-Parsing KGQA}\label{sec:path-kgqa}

\paragraph{Reasoning-on-Graphs}

Reasoning-on-Graphs (RoG)~\cite{luo2024reasoning} adopts a\index{KGQA}
planning--retrieval--reasoning pipeline. A fine-tuned model generates relation
paths grounded in the graph schema. Those paths are used to retrieve matching
instantiations. The model then reasons over the retrieved paths. Because answers
are produced from retrieved paths, RoG is faithful by construction and can
present the path as an explanation.

Two properties distinguish it from the present work. It is \emph{fine-tuned}.
This places its reported numbers outside direct comparison with prompting
methods. And its plan is generated in one shot before retrieval: if the
generated relation path does not instantiate in the graph, there is no mechanism
to revise it mid-course. The pre-generation verification question does not arise
here, because grounding is guaranteed structurally rather than checked.

RoG is nonetheless the one surveyed system this thesis shares a data
distribution with: the knowledge environment of \Cref{sec:source-data} is built
from its published subgraphs. So the comparison is taken up directly, figures
and all, in \Cref{sec:rog-comparison}.

\paragraph{StructGPT and KG-Agent}

StructGPT~\cite{jiang2023structgpt} provides a general interface for reasoning
over structured data --- knowledge graphs, tables, and databases --- through an
invoke--linearise--generate loop. Specialised interfaces extract evidence that
is linearised into text for the model. StructGPT is deliberately general rather
than graph-specific. It performs no open-ended graph search.

KG-Agent~\cite{jiang2024kgagent} is closer to the agentic setting: a
multifunctional toolbox, a knowledge-graph executor, and a dynamic memory,
driven by an instruction-tuned model that autonomously selects tools across
iterations. KG-Agent demonstrates that a small fine-tuned model with good tool
abstractions is competitive with much larger prompted ones. It plans implicitly
through tool selection rather than committing to explicit sub-objectives. It
provides no backtracking mechanism. And it performs no verification of the final
answer.

\section{Agentic Graph Exploration}\label{sec:agentic-exploration}

\subsection{Think-on-Graph}

Think-on-Graph (ToG)~\cite{sun2024think} established the template
for\index{Think-on-Graph} training-free agentic KGQA. The model performs
iterative beam search over the graph. At each step it is shown candidate
relations and entities and prunes them itself. Beam width and maximum depth are
both set to three. Exploration terminates when the model judges the retrieved
paths sufficient. At that point it answers directly.

ToG is the most directly comparable prior system to AGR --- training-free,
prompting-based, operating on Freebase, evaluated on the same benchmarks. It is
therefore the agentic baseline in \Cref{sec:baselines}. Its limitations are
instructive. There is no explicit decomposition. So multi-constraint questions
are handled implicitly. Low-scoring branches are pruned within the beam. But
there is no mechanism to return to an abandoned node once the beam has moved on.
And the transition from exploration to answering is a single model judgement:
nothing checks the answer that follows.

\subsection{Plan-on-Graph}

Plan-on-Graph (PoG)~\cite{chen2024plan} is the closest system to AGR in
ambition. It decomposes the question into sub-objectives, including constraint
conditions. It explores adaptively, with a breadth that varies by sub-objective.
And it adds Guidance, Memory, and Reflection mechanisms. Reflection decides both
whether to self-correct and which entity to backtrack to. That dual role makes
PoG the one surveyed system with genuine backtracking.

PoG's self-correction operates on the \emph{search process}: it reconsiders
where to explore. It does not decompose the draft answer into claims and check
each against the retrieved triples before responding. That difference is exactly
the gap this thesis addresses. \Cref{sec:positioning} states it precisely.

One naming hazard needs a warning here. A second, unrelated system also goes by
PoG: Paths-over-Graph~\cite{tan2024paths}. It reports considerably higher
numbers. Survey tables that cite ``PoG'' without saying which one are a known
source of muddled comparisons. This thesis therefore always writes Plan-on-Graph
or Paths-over-Graph in full.

\subsection{Generate-on-Graph and Reasoning with Trees}

Generate-on-Graph (GoG)~\cite{xu2024generate} targets the incomplete-graph
setting through a thinking--searching--generating loop: when the graph lacks a
required fact, the model generates candidate triples from its parametric
knowledge and continues. It also expands subgraphs dynamically to handle the
mediator nodes described in \Cref{sec:kgs}. ToG tends to skip these. The
trade-off is directly opposed to this thesis's: GoG deliberately admits
model-generated facts into the evidence, whereas AGR admits only graph triples.

Reasoning with Trees~\cite{shen2025reasoning} applies explicit tree search to
knowledge graph question answering. It is cited in \Cref{sec:mcts-clarification}
as an example of genuine tree-search machinery. That machinery contrasts with
the guided best-first search implemented here.

\subsection{KBQA-o1 and Genuine Tree Search}\label{sec:kbqao1}

KBQA-o1~\cite{luo2025kbqao1} is the closest system to the search strategy the
present work's proposal originally described, and therefore the most important
one to distinguish from it. It performs agentic knowledge base question
answering with full Monte Carlo Tree Search --- a policy model proposing
actions, a reward model scoring them, and rollouts whose values propagate back
through the tree --- wrapped around a ReAct-style agent that generates logical
forms step by step against a Freebase environment. The exploration additionally
generates annotations used for incremental fine-tuning. So the search improves
the model that drives it.

Two aspects bear on this thesis. First, KBQA-o1 is what the withdrawn
``MCTS-inspired'' description of \Cref{sec:mcts-clarification} would have
claimed: it has the rollouts and value backpropagation that AGR does not. Saying
so plainly is what keeps the two from being confused.

Second, its ablation speaks directly to the methodological argument of
\Cref{sec:positioning}. Removing MCTS from KBQA-o1, with every other component
left in place, drops overall GrailQA F1 from 78.5 to 48.5. That is thirty points
attributed to a single mechanism. The authors could make that attribution only
because they ablated components rather than reporting the whole assembly as one
number. This thesis applies the same instrument in \Cref{sec:ablations}. That
the two studies agree here is evidence that the instrument tells us something
real.

KBQA-o1 is not a baseline here, for three reasons. It is fine-tuned rather than
prompted. It is evaluated under a low-resource few-shot protocol. And it does
not report results on ComplexWebQuestions. So there is no way to run a
controlled comparison against it in this environment.

\subsection{Verification and Self-Correction in Language
Models}\label{sec:verification-related}

\paragraph{Chain-of-Verification}

Chain-of-Verification (CoVe)~\cite{dhuliawala2024chain} has the model draft a
response, plan verification questions about it, answer those on their own, and
then regenerate a corrected response. It establishes the core premise behind
this thesis's verification layer: break a draft into parts, check each one, and
factual errors go down. But CoVe runs its checks against the model's own
knowledge, not against an external symbolic source. Multi-agent
debate~\cite{du2024improving} gets a related effect by having several model
instances critique one another.

That distinction matters. There is evidence that language models cannot reliably
correct their own reasoning without outside
feedback~\cite{huang2024selfcorrect}. A verification step whose ground truth is
the very model that made the error is limited by construction. AGR's
verification is grounded in traversed triples. That puts the check outside the
generator.

\paragraph{Self-Correcting Graph RAG}

A self-correcting agentic graph RAG system for clinical decision support in
hepatology~\cite{hu2025selfcorrecting} shows agentic retrieval and a correction
loop working together in a high-stakes domain. It backs up the practical case
for verification, though in one specialised domain rather than on the
open-domain benchmarks used here. Reflexion-style verbal
self-improvement~\cite{shinn2023reflexion} is the general agentic pattern. The
hepatology system is a graph-specific instance of it.

\subsection{Measuring Factuality}\label{sec:factuality-related}

\paragraph{Claim-Decomposition Metrics}

FActScore~\cite{min2023factscore} evaluates long-form generation by
breaking\index{claim extraction} the text into atomic facts and checking each
one against a reference source. It reports the proportion that are supported.
This decompose-then-verify protocol is the basis for the groundedness metric in
\Cref{sec:groundedness-metric}, with the knowledge graph as the reference. The
one methodological requirement that carries over is that the judge doing the
checking must be independent of the system that wrote the text.

\paragraph{Factuality Benchmarks and Their Limits}

FactBench~\cite{bayat2025factbench} is a dynamic benchmark for factuality
evaluation in the wild. The original proposal for this thesis meant to use it
twice: as a source for the knowledge environment and as the hallucination
evaluator. That plan was dropped. The reason is methodological, not practical.
If the evaluator's facts sit inside the graph the system retrieves from, a low
hallucination rate follows automatically. It says nothing about merit. The
evaluation in \Cref{sec:setup} therefore measures claim grounding against the
graph with a human-validated judge. That keeps the evaluator's ground truth
outside the system's evidence.

\section{Comparative Summary of Agentic KGQA
Systems}\label{sec:comparison-table}

\Cref{tab:mechanisms} compares the seven closest systems on the mechanisms
relevant to this thesis. \Cref{tab:reported} records their evaluation settings
and published accuracy.

\begin{table}[!tb]
  \begin{center}
    \caption[Mechanism comparison of the seven closest prior systems]{Mechanism
    comparison of the seven closest prior systems. No system performs
    claim-level verification of the final answer against retrieved triples
    before responding.}
    \label{tab:mechanisms}
    \footnotesize
    \begin{tabular}{|l|L{38mm}|L{22mm}|L{22mm}|L{25mm}|}
      \hline
      \textbf{System} & \textbf{Exploration strategy} & \textbf{Decomposes
      query?} & \textbf{Backtracks?} & \textbf{Verifies before answering?} \\
      \hline
      ToG~\cite{sun2024think} & Training-free LLM-guided iterative beam search;
      beam width and depth both 3 & No & No --- prunes within beam, cannot
      return & No \\
      \hline
      PoG~\cite{chen2024plan} & Self-correcting adaptive planning with Guidance,
      Memory, Reflection & Yes, into sub-objectives with constraints & Yes ---
      Reflection selects the entity to return to & Partial --- corrects the
      search, not the answer \\
      \hline
      RoG~\cite{luo2024reasoning} & Planning--retrieval--reasoning over
      generated relation paths & Partial --- relation-path plans & No & Faithful
      by construction; no explicit check \\
      \hline
      KG-Agent~\cite{jiang2024kgagent} & Toolbox with KG executor and dynamic
      memory; autonomous tool selection & Implicit, via tool calls & No & No \\
      \hline
      GoG~\cite{xu2024generate} & Thinking--searching--generating; generates
      triples when the KG is incomplete & Yes & No --- expands subgraph instead
      & No \\
      \hline
      StructGPT~\cite{jiang2023structgpt} & Iterative reading-then-reasoning via
      invoke--linearise--generate interfaces & No & No & No \\
      \hline
      KBQA-o1~\cite{luo2025kbqao1} & Monte Carlo Tree Search with policy and
      reward models over a ReAct agent; fine-tuned & Implicit, stepwise logical
      form & Via MCTS tree revisiting & No \\
      \hline
      \textbf{AGR} (this work) & Guided best-first beam search with explicit
      backtracking & Yes, into ordered sub-objectives & Yes --- threshold, dead
      end, or evaluator veto & \textbf{Yes --- claim-level, pre-generation} \\
      \hline
    \end{tabular}
  \end{center}
\end{table}

\begin{table}[!tb]
  \begin{center}
    \caption[Reported evaluation settings and Hits@1 as published]{Reported
    evaluation settings and Hits@1 as published. These figures are \emph{not}
    directly comparable to one another: backbones differ, some papers evaluate
    on sampled subsets, and the backbone accounts for much of the variation.}
    \label{tab:reported}
    \footnotesize
    \begin{tabular}{|l|l|L{45mm}|c|c|}
      \hline
      \textbf{System} & \textbf{Backbone} & \textbf{Datasets} & \textbf{WebQSP}
      & \textbf{CWQ} \\
      \hline
      ToG & GPT-3.5 & CWQ, WebQSP, GrailQA~\cite{gu2021beyond}, QALD10-en,
      SimpleQuestions, WebQuestions, T-REx, Zero-Shot RE, Creak & 76.2 & 58.9 \\
      \cline{2-5}
       & GPT-4 & & 82.6 & 69.5 \\
      \hline
      PoG & GPT-3.5 & CWQ, WebQSP, GrailQA & 82.0 & 63.2 \\
      \cline{2-5}
       & GPT-4 & & 87.3 & 75.0 \\
      \hline
      RoG & Fine-tuned LLaMA2-7B & WebQSP, CWQ & 85.7 & 62.6 \\
      \hline
      KG-Agent & Fine-tuned 7B & WebQSP, CWQ, plus other KBs & 83.3 & 72.2 \\
      \hline
      GoG & GPT-3.5 & WebQSP, CWQ (complete and incomplete KG) & 78.7 & 55.7 \\
      \cline{2-5}
       & GPT-4 & & 84.4 & 75.2 \\
      \hline
      StructGPT & GPT-3.5 & WebQSP, MetaQA, TableQA, Text-to-SQL & 72.6 & --- \\
      \hline
      KBQA-o1 & Llama-3.1-8B & GrailQA, WebQSP, GraphQ; 100-shot low-resource
      protocol & 59.8$^{\dagger}$ & --- \\
      \cline{2-5}
       & Llama-3.3-70B & & 67.0$^{\dagger}$ & --- \\
      \hline
    \end{tabular}

    \vspace{1mm}
    {\footnotesize $^{\dagger}$ KBQA-o1 reports \textbf{F1}, not Hits@1, and
    evaluates under a 100-shot low-resource protocol. These cells are therefore
    \emph{not} comparable with the rest of the column and are shown only to
    situate the system. It does not evaluate CWQ.

    \vspace{0.7mm}
    Every figure is the one its own authors published. That this is not the same
    as a comparison is visible inside the table: PoG's paper reproduces
    Think-on-Graph and measures it at $57.1$ and $67.6$ on CWQ, against the
    $58.9$ and $69.5$ Think-on-Graph reports for itself. Two published papers
    disagree by $1.8$ and $1.9$ points on the same system and benchmark. That
    gap is the order of many claimed improvements in this literature. It is also
    the reason \Cref{sec:setup} reimplements its baseline rather than citing a
    number.}
  \end{center}
\end{table}

Three caveats attach to \Cref{tab:reported} and are carried into every
comparison in this thesis.

First, the figures are drawn from the original publications and reflect
different backbones. RoG and KG-Agent are fine-tuned. Their numbers are not
comparable with prompting methods at all. Among the prompting methods, the
backbone accounts for much of the spread: the same methods drop sharply when run
on smaller open-weight models. That sensitivity is the direct motivation for
holding the backbone constant across every system in \Cref{sec:setup}. The
literature does not hold it constant. So a cross-paper table cannot settle which
architecture is better.

Second, evaluation protocols differ. Some papers evaluate on sampled test
subsets. And answer-matching conventions are not uniform across implementations.

Third, and most concretely, two papers disagree about a third. Plan-on-Graph's
own comparison table measures Think-on-Graph at $57.1$ and $67.6$ on CWQ, while
Think-on-Graph reports $58.9$ and $69.5$ for itself. Neither is wrong: one is an
original result and the other a reimplementation. This thesis is in exactly that
situation with respect to its own agentic baseline (\Cref{sec:baseline-tog}).
The gap is $1.8$ and $1.9$ points, the same order as many improvements claimed
in this literature. A cross-paper table cannot distinguish a real advance from a
reproduction difference of that size. That is why no claim in this thesis rests
on one.

\subsection{Positioning of This Work}\label{sec:positioning}

\Cref{tab:mechanisms} makes the gap explicit. Decomposition is well established:
PoG and GoG both do it. Backtracking exists, in PoG. Faithful grounding by
construction exists, in RoG. But in the final column, \emph{no surveyed system
decomposes its draft answer into atomic claims and checks each against the
retrieved triples before responding}. PoG's Reflection corrects the search while
it is running. RoG's answers are grounded because of how they are produced, not
because anything examined them. The slot is open. \Cref{sec:contribution} states
what this thesis puts in it: the Structural Verification Layer
(\Cref{sec:verification-layer}) together with a measurement of what it buys, a
component ablation pricing each agentic mechanism in accuracy and tokens
(\Cref{sec:ablations}), and the controlled agentic comparison itself
(\Cref{sec:findings}).

Two features of that positioning belong here rather than there. Both are claims
about the literature. The first is that these mechanisms are typically reported
as bundles. So which of them earns its cost is not established. That is the gap
the ablation fills, and where the planner's stratum dependence appears. The
second is that the thesis does not depend on beating prior systems on raw
Hits@1. The comparison that matters is against the same architecture with the
verification layer removed, on the same graph, with the same backbone. That is
why \Cref{sec:ablations} rather than \Cref{tab:reported} carries the argument.
Where AGR is compared against a prior system, that comparison is run under
controlled conditions in this thesis's own environment rather than quoted across
papers.

\endinput
